\section{Introduction}

Recently, large language models (LLMs)~\cite{zhao2023survey} have demonstrated exceptional capabilities across a variety of tasks within the field of natural language processing. However, \emph{hallucination}~\cite{ji2023survey, rawte2023survey, zhang2023siren, huang2023survey, ye2023cognitive} in text generated by LLMs remains an underlying concern, impeding the application of LLMs in real-world scenarios~\cite{kaddour2023challenges}.
\citet{xu2024hallucination} have indicated that, despite the existence of some effective hallucination mitigation strategies, the occurrence of hallucinations in LLMs is inevitable. Therefore, reliable and effective hallucination detection methods are necessary and urgent.

% Most existing hallucination detection methods depend heavily on powerful LLMs such as GPT-4 and can be broadly categorized into two primary approaches. 
Existing hallucination detection methods can be roughly categorized into two primary approaches. 
One line of work relies on the internal knowledge of LLMs to directly identify hallucinations via prompts~\cite{li2023halueval, lei2023chain} or evaluating the semantic consistency among multiple responses to the same question generated by LLMs~\cite{manakul2023selfcheckgpt}. 
% and consistency~\cite{wang2022self} to detect hallucinations. This involves asking LLMs to identify hallucinations in the text~\cite{li2023halueval, lei2023chain} or generating multiple responses to the same question and comparing their consistency to detect hallucinations~\cite{manakul2023selfcheckgpt}. 
% However, these methods rely too heavily on a single judgment by the model, constrained not only by LLMs' internal knowledge boundaries but also by its ability to utilize that knowledge. For texts beyond the model's scope of knowledge, they may not accurately identify hallucinations.
However, these methods are usually constrained not only by LLMs' internal knowledge but also by their abilities to utilize knowledge.
Another line of work extends the detection ability of LLMs by employing external tools (\eg search engine) to obtain supporting evidence for hallucination detection~\cite{chern2023factool, wei2024long, min2023factscore}. However, these approaches mostly depend on closed-source powerful LLMs such as GPT-4. Moreover, their detection process is usually pre-determined by human, making it difficult for the model to autonomously and effectively execute the hallucination detection. 

To address these issues, in this paper, we propose the \textbf{HaluAgent}, an autonomous LLM-based agent framework for hallucination detection. This agent is based on smaller open-source models (\ie Baichuan2-Chat 7B and 13B~\cite{yang2023baichuan}) and capable of \emph{bilingual} hallucination detection in Chinese and English. The motivations are twofold: (1) designing autonomous detection agents that can actively make decisions and judgements, without human
assistance; (2) enabling relatively smaller models to effectively perform complex detection, without reliance on close-sourced LLM APIs. To achieve this, we make three major technical contributions. First, we extend the LLM's capability to detect a broader range of hallucination forms such as text, code, math expression, or their combination by curating a multi-functional toolbox, in contrast to previous work only focused on textual hallucination or limited tools~\cite{manakul2023selfcheckgpt,min2023factscore}. Second, we design a fine-grained three-stage detection framework along with memory mechanism, including sentence segmentation, tool selection and verification, and reflection. Third, we leverage existing hallucination datasets to synthesize detection trajectories for fine-tuning the LLM, where we first employ GPT-4 to execute the above three stages until obtaining detection results consistent with the ground-truth label and then synthesize the instruction data. We compare HaluAgent and previous work in Table~\ref{tab: detection-methods}.

To verify the effectiveness, We evaluate HaluAgent on both in-domain and out-of-domain datasets at response- and sentence-level granularities. After fine-tuning with only 2K trajectories, the detection performance of HaluAgent has been improved significantly (\eg overall accuracy increases from 46.44\% to 79.70\% in four in-domain datasets and from 49.50\% to 78.43\% in two out-of-domain datasets), reaching a level comparable to or even higher than GPT-4 without tool enhancement. In the sentence-level detection experiments, HaluAgent also achieved substantial improvements, particularly with F1 scores on the math and science datasets increasing from 19.51\% to 68.80\% and from 17.54\% to 94.16\%, respectively.


% In open-ended real-world interactions between users and LLMs, texts requiring hallucination detection are often diverse in type and language. Therefore, hallucination detection is a highly complex and challenging task that demands LLMs to flexibly leverage internal knowledge or appropriate tools to achieve accurate detection. Even powerful LLMs like GPT-4 find this task challenging, and the challenge escalates for smaller open-source LLMs (\eg 7B models).
% Thus, a flexible hallucination detector capable of handling this diversity is needed.


% To overcome the limitations of existing methods, we propose treating the hallucination detection task as an agent task. 
% On one hand, the agent can invoke various tools, breaking the limitation of relying on world knowledge inherent in LLMs for hallucination detection. 
% On the other hand, the agent possesses planning and reasoning abilities and can adjust its actions based on the environment, allowing it to flexibly detect hallucinations in various types of text. 
% With these capabilities, small open-source models (\eg 7B models) can effectively perform hallucination detection tasks without relying on closed-source LLM APIs.
% To overcome the limitations of existing methods, we propose treating hallucination detection as an agent task. We assert that by equipping small open-source models with planning, reasoning, and tool utilization capabilities of agents, they can effectively perform hallucination detection.
% To achieve this, we introduce HaluAgent, a flexible and universal bilingual hallucination detection agent capable of handling various types of text. Our work includes three major technical contributions.
% Firstly, we design an automated agent framework equipped with a flexible toolbox specifically for hallucination detection tasks. 
% The input text is first segmented into sentences, and each sentence is verified using the appropriate tool. Detection results of each sentence are stored in the agent's memory and reviewed to produce a final outcome. If hallucinations are detected, the specific sentences and supporting evidence provided by the tools will be output. The tool invocation and reflection mechanisms allow LLMs to integrate external knowledge and reassess their judgments, while the agent framework enables LLMs to progressively plan and reason on diverse texts. 
% These designs address the limitations of LLMs in hallucination detection tasks.
% We present the similarities and differences between HaluAgent and other detection methods in Table~\ref{tab: detection-methods}.


\begin{table}[tb]
\centering
\small
\renewcommand\arraystretch{1.15}
\setlength\tabcolsep{2.0pt}
\resizebox{0.49\textwidth}{!}{
\begin{tabular}{l|ccccc}
\toprule
\textbf{Methods}  & \textbf{\makecell{Base\\Model}}& \textbf{\makecell{Task\\ Agnostic}} & \textbf{\makecell{Tool\\ Usage}} & \textbf{\makecell{Fine\\ Grained}} & \textbf{\makecell{Extensi-\\bility}}\\
\midrule
SelfCheckGPT & ChatGPT & \usym{2717} & \usym{2717} & \usym{2717} &  \usym{2717}\\
SAFE & GPT-4 & \usym{2717} & \usym{2713} & \usym{2713} & \usym{2717} \\
FacTool & GPT-4 & \usym{2717}  & \usym{2713} & \usym{2713} & \usym{2717}\\
HaluAgent & Baichuan2-Chat  & \usym{2713}  & \usym{2713} & \usym{2713} & \usym{2713} \\
\bottomrule
\end{tabular}
}
\caption{Comparison of different methods. \textbf{{Task Agnostic}} means whether the method is designed to specific tasks; \textbf{{Fine Grained}} describes whether providing detailed hallucination sentences; \textbf{{Extensibility}} means whether the method can extend to more tasks and tools.}
\label{tab: detection-methods}
% \vspace{-0.2cm}
\end{table}


% Secondly, we leverage existing datasets and ChatGPT to construct a bilingual hallucination detection dataset, incorporating various tasks such as knowledge-based QA, math problem solving, code generation, and conditional text generation. Following our proposed framework, we employ GPT-4 to generate hallucination detection trajectories in ReAct~\cite{yao2023react} format with \textit{thought}, \textit{action}, and \textit{observation}. To ensure quality, we retain only the trajectories that successfully executed the detection process, resulting in a high-quality hallucination detection agent trajectory dataset.

% Thirdly, we fine-tune Baichuan2-Chat (7B and 13B) on the synthesized trajectories and evaluate the resulting models on both in-domain and out-of-domain tasks. Unlike other hallucination detection methods, HaluAgent is capable of handling hallucination detection for various types of text in both Chinese and English. After trajectory fine-tuning, the hallucination detection capability of small open-source models has been significantly improved, reaching a level comparable to GPT-4 without tool enhancements.





